% ============================================================
% Section 6: Backup Slides (Optional)
% ============================================================

\begin{frame}{Backup: Adaptive Bandwidth}
    \textbf{Problem:} Fixed bandwidth may not fit all regions

    \vspace{0.3cm}

    \textbf{Solution:} Adaptive Mean Shift
    \[
    h_i = h_0 \cdot \left(\frac{\hat{f}(\mathbf{x}_i)}{\text{geom\_mean}(\hat{f})}\right)^{-\alpha}
    \]

    \begin{itemize}
        \item Smaller bandwidth in dense regions
        \item Larger bandwidth in sparse regions
        \item Parameter $\alpha$ controls adaptation strength
    \end{itemize}

    \vspace{0.5cm}

    \textbf{Advantages:}
    \begin{itemize}
        \item Better handling of varying densities
        \item More robust to bandwidth choice
    \end{itemize}

    \textbf{Disadvantages:}
    \begin{itemize}
        \item Additional parameter to tune ($\alpha$)
        \item Increased computational cost
    \end{itemize}
\end{frame}

\begin{frame}{Backup: Implementation Tips}
    \textbf{Performance Optimizations:}
    \begin{enumerate}
        \item \textbf{Spatial indexing:} Use KD-trees for neighbor search
        \begin{itemize}
            \item Reduces complexity from $O(n^2)$ to $O(n \log n)$
        \end{itemize}

        \item \textbf{Binning:} Discretize space, compute shift on bins
        \begin{itemize}
            \item Trade accuracy for speed
        \end{itemize}

        \item \textbf{Early stopping:} Merge nearby modes
        \begin{itemize}
            \item Use threshold to merge modes within distance $\epsilon$
        \end{itemize}

        \item \textbf{Sampling:} Use subset for large datasets
        \begin{itemize}
            \item Cluster sample, then assign remaining points
        \end{itemize}
    \end{enumerate}

    \vspace{0.5cm}

    \textbf{R Implementation:}
    \begin{itemize}
        \item \texttt{MeanShift} package
        \item \texttt{meanShiftR} (faster, C++ backend)
        \item Custom implementation for learning
    \end{itemize}
\end{frame}

\begin{frame}{Backup: Comparison with Other Methods}
    \begin{table}
        \centering
        \small
        \begin{tabular}{lcccc}
            \toprule
            \textbf{Algorithm} & \textbf{Clusters} & \textbf{Shape} & \textbf{Speed} & \textbf{Parameters} \\
            \midrule
            K-means & Predefined & Spherical & Fast & $k$ \\
            DBSCAN & Automatic & Arbitrary & Medium & $\epsilon$, minPts \\
            Mean Shift & Automatic & Arbitrary & Slow & $h$ \\
            Hierarchical & Flexible & Arbitrary & Slow & Linkage, $k$ \\
            Gaussian Mixture & Predefined & Elliptical & Medium & $k$ \\
            \bottomrule
        \end{tabular}
    \end{table}

    \vspace{0.5cm}

    \textbf{Key Takeaway:}
    \begin{itemize}
        \item Mean Shift offers good balance of flexibility and simplicity
        \item Best when cluster count is unknown and shapes are arbitrary
        \item Speed is main limitation for large datasets
    \end{itemize}
\end{frame}

\begin{frame}{Backup: Mathematical Details}
    \textbf{Convergence Proof Sketch:}

    \begin{enumerate}
        \item Mean shift is gradient ascent on kernel density estimate
        \item Gradient of KDE:
        \[
        \nabla \hat{f}(\mathbf{x}) = \frac{1}{nh^{d+2}} \sum_{i=1}^{n} (\mathbf{x} - \mathbf{x}_i) K'\left(\frac{\|\mathbf{x} - \mathbf{x}_i\|}{h}\right)
        \]

        \item Mean shift direction aligns with gradient
        \item Converges to local mode (stationary point where $\nabla \hat{f} = 0$)
    \end{enumerate}

    \vspace{0.5cm}

    \textbf{Convergence Rate:}
    \begin{itemize}
        \item Linear convergence near modes
        \item Typically 10-50 iterations in practice
        \item Depends on bandwidth and data structure
    \end{itemize}
\end{frame}
